\documentclass[]{article}

\usepackage{amsmath}
\usepackage{amsfonts}
\usepackage{amsthm}
\usepackage{amssymb}
\usepackage{mathrsfs}
\usepackage{stmaryrd}

\newcommand{\Q}{\mathbb{Q}}
\newcommand{\N}{\mathbb{N}}
\newcommand{\Z}{\mathbb{Z}}
\newcommand{\R}{\mathbb{R}}
\newcommand{\Primes}{\mathbb{P}}
\newcommand{\st}{\text{ s.t. }}
\newcommand{\txtand}{\text{ and }}
\newcommand{\txtor}{\text{ or }}
\newcommand{\lxor}{\veebar}

%opening
\title{Notes on Linear Algebra}
\author{DSBA Mathematics Refresher 2025}
\date{}

\begin{document}
	
	\maketitle
	
	\begin{abstract}
		
	\end{abstract}
	
	
	\section{Linear Independence Property}
	
	Given a set of vectors $\{\mathbf{v}_1, \mathbf{v}_2, \dots, \mathbf{v}_n\}$ in a vector space $V$, the set is said to be \textbf{linearly independent} if the only solution to the equation
	$$
	c_1\mathbf{v}_1 + c_2\mathbf{v}_2 + \dots + c_n\mathbf{v}_n = \mathbf{0}
	$$
	is 
	$$
	c_1 = c_2 = \dots = c_n = 0.
	$$\\
	In other words, none of the vectors in the set can be written as a linear combination of the others.
	
	\noindent \textbf{Example:}
	Consider the vectors
	$\mathbf{v}_1 = \begin{pmatrix} 3 \\ 2 \\ 4 \\ -1 \\ 2 \end{pmatrix}$,
	$\mathbf{v}_2 = \begin{pmatrix} -2 \\ 5 \\ 0 \\ 1 \\ 1 \end{pmatrix}$
	and
	$\mathbf{v}_3 = \begin{pmatrix} 1 \\ 0 \\ 5 \\ 3 \\ -4 \end{pmatrix}$
	in $\mathbb{R}^5$.
	These vectors are linearly independent because the only solution to
	$$
	c_1v_1 + c_2v_2 + c_3v_3 = \mathbf{0}
	$$
	is $c_1 = c_2 = c_3 = 0$.
	Indeed, writing the equation component by component gives the system
	$$
	\begin{cases}
		3c_1 - 2c_2 + c_3 &= 0 \\
		2c_1 + 5c_2 &= 0 \\
		4c_1 + 5c_3 &= 0 \\
		-c_1 + c_2 + 3c_3 &= 0 \\
		2c_1 + c_2 - 4c_3 &= 0
	\end{cases}
	$$
	The second line gives $c_2 = -\frac{2}{5}c_1$ and the third gives $c_3 = -\frac{4}{5}c_1$;
	substituting both into the first line:
	$$
	3c_1 + \frac{4}{5}c_1 - \frac{4}{5}c_1 = 3c_1 = 0
	\qquad \Longrightarrow \qquad
	c_1 = 0, \quad \text{and hence} \quad c_2 = c_3 = 0
	$$

	\noindent \textbf{Counter-example:}
	The vectors
	$\mathbf{w}_1 = \begin{pmatrix} 1 \\ 0 \\ 1 \end{pmatrix}$,
	$\mathbf{w}_2 = \begin{pmatrix} 0 \\ 1 \\ 1 \end{pmatrix}$
	and
	$\mathbf{w}_3 = \begin{pmatrix} 1 \\ 1 \\ 2 \end{pmatrix}$
	in $\mathbb{R}^3$ are \textbf{not} linearly independent, because
	$$
	\mathbf{w}_1 + \mathbf{w}_2 - \mathbf{w}_3 = \mathbf{0}
	$$
	is a non-trivial linear combination giving $\mathbf{0}$ (here $c_1 = 1$, $c_2 = 1$, $c_3 = -1$).
	Equivalently, $\mathbf{w}_3 = \mathbf{w}_1 + \mathbf{w}_2$: the third vector carries no information that the first two do not already carry.
	
	\section{Spanning Property}
	
	A set of vectors $\{\mathbf{v}_1, \mathbf{v}_2, \dots, \mathbf{v}_n\}$ in a vector space $V$ is said to \textbf{span} $V$ if every vector in $V$ can be expressed as a linear combination of these vectors.
	Mathematically, the set spans $V$ if for every vector $\mathbf{v} \in V$, there exist scalars $c_1, c_2, \dots, c_n$ such that
	$$
	\mathbf{v} = c_1\mathbf{v}_1 + c_2\mathbf{v}_2 + \dots + c_n\mathbf{v}_n.
	$$
	
	\noindent \textbf{Example:}
	In $\mathbb{R}^2$, the vectors
	$\mathbf{v}_1 = \begin{pmatrix} 1 \\ 1 \end{pmatrix}$,
	$\mathbf{v}_2 = \begin{pmatrix} -1 \\ 2 \end{pmatrix}$
	and
	$\mathbf{v}_3 = \begin{pmatrix} 1 \\ -1 \end{pmatrix}$
	span $\mathbb{R}^2$ because any vector
	$\mathbf{v} = \begin{pmatrix} x \\ y \end{pmatrix}$
	can be written as
	$$
	\mathbf{v} = \frac{2x+y}{3}\mathbf{v}_1 + \frac{y-x}{3}\mathbf{v}_2 + 0 \cdot \mathbf{v}_3.
	$$
	Indeed,
	$$
	\frac{2x+y}{3}\begin{pmatrix} 1 \\ 1 \end{pmatrix} + \frac{y-x}{3}\begin{pmatrix} -1 \\ 2 \end{pmatrix}
	= \begin{pmatrix} \frac{2x+y-y+x}{3} \\ \frac{2x+y+2y-2x}{3} \end{pmatrix}
	= \begin{pmatrix} x \\ y \end{pmatrix}.
	$$
	Note that such a decomposition is \textbf{not unique} here: the three vectors are redundant (one checks that $\mathbf{v}_3 = \frac{1}{3}\mathbf{v}_1 - \frac{2}{3}\mathbf{v}_2$), and we could equally have written $\mathbf{v} = \frac{x+y}{2}\mathbf{v}_1 + 0 \cdot \mathbf{v}_2 + \frac{x-y}{2}\mathbf{v}_3$.
	Spanning only requires that \textit{at least one} decomposition exists.

	\noindent \textbf{Counter-example:}
	In $\mathbb{R}^3$, the vectors
	$\mathbf{w}_1 = \begin{pmatrix} 1 \\ 0 \\ 0 \end{pmatrix}$
	and
	$\mathbf{w}_2 = \begin{pmatrix} 0 \\ 1 \\ 0 \end{pmatrix}$
	do \textbf{not} span $\mathbb{R}^3$: every linear combination $c_1\mathbf{w}_1 + c_2\mathbf{w}_2 = \begin{pmatrix} c_1 \\ c_2 \\ 0 \end{pmatrix}$ has a third component equal to $0$, so for instance $\begin{pmatrix} 0 \\ 0 \\ 1 \end{pmatrix}$ cannot be reached.
	They span the plane $\{z = 0\}$, which is a strict subspace of $\mathbb{R}^3$.
	
	\section{Basis}
	
	A set of vectors $\{\mathbf{v}_1, \mathbf{v}_2, \dots, \mathbf{v}_n\}$ in a vector space $V$ is called a \textbf{basis} for $V$ if:
	\begin{itemize}
		\item The set is linearly independent.
		\item The set spans the vector space $V$.
	\end{itemize}
	
	The number of vectors in a basis is called the \textbf{dimension} of the vector space.
	
	\noindent \textbf{Example:}
	The vectors $\mathbf{v}_1 = \begin{pmatrix} 1 \\ 0 \end{pmatrix}$ and $\mathbf{v}_2 = \begin{pmatrix} 0 \\ 1 \end{pmatrix}$ form a basis for $\mathbb{R}^2$.
	The dimension of $\mathbb{R}^2$ is 2.

	\noindent \textbf{Example (a basis is not unique):}
	The vectors $\mathbf{u}_1 = \begin{pmatrix} 1 \\ 1 \end{pmatrix}$ and $\mathbf{u}_2 = \begin{pmatrix} 1 \\ -1 \end{pmatrix}$ also form a basis for $\mathbb{R}^2$:
	\begin{itemize}
		\item They are linearly independent: $c_1\mathbf{u}_1 + c_2\mathbf{u}_2 = \mathbf{0}$ gives $c_1 + c_2 = 0$ and $c_1 - c_2 = 0$, so $c_1 = c_2 = 0$.
		\item They span $\mathbb{R}^2$: any $\begin{pmatrix} x \\ y \end{pmatrix} = \frac{x+y}{2}\mathbf{u}_1 + \frac{x-y}{2}\mathbf{u}_2$.
	\end{itemize}
	A vector space has many different bases, but they all have the \textit{same number of vectors}, which is why the dimension is well defined.

	\noindent \textbf{Counter-examples:}
	\begin{itemize}
		\item The three vectors $\mathbf{v}_1, \mathbf{v}_2, \mathbf{v}_3$ of the previous section span $\mathbb{R}^2$, but they are \textit{not} a basis: $3$ vectors in a $2$-dimensional space can never be linearly independent.
		\item The two vectors $\mathbf{w}_1, \mathbf{w}_2$ of the previous section are linearly independent in $\mathbb{R}^3$, but they are \textit{not} a basis of $\mathbb{R}^3$ either: they do not span it.
	\end{itemize}
	In general, in a space of dimension $n$, a basis has \textit{exactly} $n$ vectors: fewer cannot span, more cannot be independent.
	
	\section{(Reduced) Row Echelon Form}
	
	A matrix is in \textbf{row echelon form} if it satisfies the following conditions:
	\begin{enumerate}
		\item All non-zero rows are above any rows of all zeros.
		\item The leading entry of each non-zero row after the first occurs to the right of the leading entry of the previous row.
		\item The leading entry in any non-zero row is 1 (this condition is for the \textbf{reduced} row echelon form).
		\item The leading 1 is the only non-zero entry in its column (for \textbf{reduced} row echelon form).
	\end{enumerate}
	
	\textbf{Example:} The matrix
	$$
	\begin{pmatrix}
		1 & 2 & 3 \\
		0 & 1 & 4 \\
		0 & 0 & 1
	\end{pmatrix}
	$$
	is in row echelon form.
	It is \textit{not} in \textbf{reduced} row echelon form, because the leading $1$'s are not the only non-zero entries of their columns.
	Clearing the entries above the pivots, with $L_2 \leftarrow L_2 - 4L_3$, then $L_1 \leftarrow L_1 - 3L_3$ and finally $L_1 \leftarrow L_1 - 2L_2$, gives the reduced row echelon form:
	$$
	\begin{pmatrix}
		1 & 0 & 0 \\
		0 & 1 & 0 \\
		0 & 0 & 1
	\end{pmatrix}
	$$

	\noindent \textbf{Example (a full reduction):}
	Let us reduce
	$$
	M = \begin{pmatrix}
		1 & 2 & -1 & 3 \\
		2 & 4 & 1 & 0 \\
		3 & 6 & 2 & 1
	\end{pmatrix}
	$$
	First, use the pivot in position $(1,1)$ to clear the entries below it, with $L_2 \leftarrow L_2 - 2L_1$ and $L_3 \leftarrow L_3 - 3L_1$:
	$$
	\begin{pmatrix}
		1 & 2 & -1 & 3 \\
		0 & 0 & 3 & -6 \\
		0 & 0 & 5 & -8
	\end{pmatrix}
	$$
	The second column has no pivot (it is a \textbf{free} column), so we move to the third one.
	With $L_2 \leftarrow \frac{1}{3} L_2$, then $L_3 \leftarrow L_3 - 5L_2$ and $L_3 \leftarrow \frac{1}{2}L_3$:
	$$
	\begin{pmatrix}
		1 & 2 & -1 & 3 \\
		0 & 0 & 1 & -2 \\
		0 & 0 & 0 & 2
	\end{pmatrix}
	\qquad \longrightarrow \qquad
	\begin{pmatrix}
		1 & 2 & -1 & 3 \\
		0 & 0 & 1 & -2 \\
		0 & 0 & 0 & 1
	\end{pmatrix}
	$$
	This is a row echelon form.
	To make it \textit{reduced}, we clear the entries above each pivot, with $L_2 \leftarrow L_2 + 2L_3$, then $L_1 \leftarrow L_1 + L_2$ and $L_1 \leftarrow L_1 - 3L_3$:
	$$
	\begin{pmatrix}
		1 & 2 & 0 & 0 \\
		0 & 0 & 1 & 0 \\
		0 & 0 & 0 & 1
	\end{pmatrix}
	$$
	which is the reduced row echelon form of $M$ (with pivots in columns 1, 3 and 4).

	\noindent \textbf{Example (solving a system):}
	This is what row reduction is mostly used for.
	To solve
	$$
	\begin{cases}
		x + 2y + z &= 8 \\
		2x + y - z &= 1 \\
		3x - y + 2z &= 7
	\end{cases}
	$$
	we write the \textbf{augmented matrix} of the system, and reduce it.
	With $L_2 \leftarrow L_2 - 2L_1$ and $L_3 \leftarrow L_3 - 3L_1$:
	$$
	\left(\begin{array}{ccc|c}
		1 & 2 & 1 & 8 \\
		2 & 1 & -1 & 1 \\
		3 & -1 & 2 & 7
	\end{array}\right)
	\qquad \longrightarrow \qquad
	\left(\begin{array}{ccc|c}
		1 & 2 & 1 & 8 \\
		0 & -3 & -3 & -15 \\
		0 & -7 & -1 & -17
	\end{array}\right)
	$$
	With $L_2 \leftarrow -\frac{1}{3}L_2$, then $L_3 \leftarrow L_3 + 7L_2$ and $L_3 \leftarrow \frac{1}{6}L_3$:
	$$
	\left(\begin{array}{ccc|c}
		1 & 2 & 1 & 8 \\
		0 & 1 & 1 & 5 \\
		0 & -7 & -1 & -17
	\end{array}\right)
	\qquad \longrightarrow \qquad
	\left(\begin{array}{ccc|c}
		1 & 2 & 1 & 8 \\
		0 & 1 & 1 & 5 \\
		0 & 0 & 1 & 3
	\end{array}\right)
	$$
	The last row reads $z = 3$; substituting back into the second row, $y + 3 = 5$, so $y = 2$; and into the first, $x + 4 + 3 = 8$, so $x = 1$.
	This last stage is called \textbf{back-substitution}.
	$$
	x = 1, \qquad y = 2, \qquad z = 3
	$$
	\textit{Check:} $1 + 4 + 3 = 8$ \checkmark, \ $2 + 2 - 3 = 1$ \checkmark, \ $3 - 2 + 6 = 7$. \checkmark
	
	\section{Inverting a Matrix}
	
	The \textbf{inverse} of a square matrix $A$ is denoted by $A^{-1}$ and is defined as the matrix that satisfies
	$$
	AA^{-1} = A^{-1}A = I,
	$$
	where $I$ is the identity matrix.
	
	\textbf{Steps to find the inverse of a matrix:}
	\begin{enumerate}
		\item Form the augmented matrix $[A | I]$.
		\item Use row operations to convert $A$ into the identity matrix.
		\item The matrix that results from the identity matrix on the left side is $A^{-1}$ on the right side.
	\end{enumerate}
	
	\noindent \textbf{Example:} For the matrix
	$$
	A = \begin{pmatrix} 1 & 2 \\ 3 & 4 \end{pmatrix},
	$$
	the augmented matrix is
	$$
	[A | I] = \begin{pmatrix}
		1 & 2 & \mid & 1 & 0\\
		3 & 4 & \mid & 0 & 1
	\end{pmatrix}
	$$
	after reducing to row echelon form:
	$$
	[I | A^{-1}] = \begin{pmatrix}
		1 & 0 & \mid & -2 & 1\\
		0 & 1 & \mid & 1.5 & -0.5
	\end{pmatrix}
	$$
	hence, the inverse is
	$$
	A^{-1} = \begin{pmatrix}
		-2 & 1\\
		1.5 & -0.5
	\end{pmatrix}
	$$
	In detail, the row operations are $L_2 \leftarrow L_2 - 3L_1$, then $L_2 \leftarrow -\frac{1}{2}L_2$, and finally $L_1 \leftarrow L_1 - 2L_2$:
	$$
	\left(\begin{array}{cc|cc}
		1 & 2 & 1 & 0\\
		0 & -2 & -3 & 1
	\end{array}\right)
	\longrightarrow
	\left(\begin{array}{cc|cc}
		1 & 2 & 1 & 0\\
		0 & 1 & 1.5 & -0.5
	\end{array}\right)
	\longrightarrow
	\left(\begin{array}{cc|cc}
		1 & 0 & -2 & 1\\
		0 & 1 & 1.5 & -0.5
	\end{array}\right)
	$$
	It is always worth checking the result:
	$$
	A A^{-1} = \begin{pmatrix} 1 & 2 \\ 3 & 4 \end{pmatrix}\begin{pmatrix} -2 & 1\\ 1.5 & -0.5 \end{pmatrix}
	= \begin{pmatrix} -2+3 & 1-1 \\ -6+6 & 3-2 \end{pmatrix}
	= \begin{pmatrix} 1 & 0 \\ 0 & 1 \end{pmatrix} = I \ \checkmark
	$$

	\noindent \textbf{Example ($3 \times 3$):} For
	$$
	N = \begin{pmatrix} 1 & 2 & 3 \\ 0 & 1 & 4 \\ 5 & 6 & 0 \end{pmatrix},
	$$
	we start from $[N | I]$ and apply $L_3 \leftarrow L_3 - 5L_1$, then $L_3 \leftarrow L_3 + 4L_2$:
	$$
	\left(\begin{array}{ccc|ccc}
		1 & 2 & 3 & 1 & 0 & 0\\
		0 & 1 & 4 & 0 & 1 & 0\\
		5 & 6 & 0 & 0 & 0 & 1
	\end{array}\right)
	\longrightarrow
	\left(\begin{array}{ccc|ccc}
		1 & 2 & 3 & 1 & 0 & 0\\
		0 & 1 & 4 & 0 & 1 & 0\\
		0 & -4 & -15 & -5 & 0 & 1
	\end{array}\right)
	\longrightarrow
	\left(\begin{array}{ccc|ccc}
		1 & 2 & 3 & 1 & 0 & 0\\
		0 & 1 & 4 & 0 & 1 & 0\\
		0 & 0 & 1 & -5 & 4 & 1
	\end{array}\right)
	$$
	The left block is now in row echelon form; we finish with $L_1 \leftarrow L_1 - 2L_2$, then $L_2 \leftarrow L_2 - 4L_3$ and $L_1 \leftarrow L_1 + 5L_3$:
	$$
	\left(\begin{array}{ccc|ccc}
		1 & 0 & -5 & 1 & -2 & 0\\
		0 & 1 & 4 & 0 & 1 & 0\\
		0 & 0 & 1 & -5 & 4 & 1
	\end{array}\right)
	\longrightarrow
	\left(\begin{array}{ccc|ccc}
		1 & 0 & 0 & -24 & 18 & 5\\
		0 & 1 & 0 & 20 & -15 & -4\\
		0 & 0 & 1 & -5 & 4 & 1
	\end{array}\right)
	$$
	hence
	$$
	N^{-1} = \begin{pmatrix} -24 & 18 & 5 \\ 20 & -15 & -4 \\ -5 & 4 & 1 \end{pmatrix}
	$$
	(one checks that $N N^{-1} = I$).
	
	\section{Determinant}
	
	The \textbf{determinant} of a square matrix $A$, denoted by $\text{det}(A)$, is a scalar value that can be computed from the elements of the matrix. The determinant has important properties and applications, including:
	\begin{itemize}
		\item A matrix is invertible if and only if its determinant is non-zero.
		\item The determinant of a product of matrices is the product of their determinants: $\text{det}(AB) = \text{det}(A)\text{det}(B)$.
	\end{itemize}
	
	For a $2 \times 2$ matrix $A = \begin{pmatrix} a & b \\ c & d \end{pmatrix}$, the determinant is calculated as:
	$$
	\text{det}(A) = ad - bc.
	$$
	
	For a $3 \times 3$ matrix $A = \begin{pmatrix} a & b & c \\ d & e & f \\ g & h & i \end{pmatrix} $, the determinant is calculated as:
	$$
	\text{det}(A) =
	  a \begin{vmatrix} e & f \\ h & i \end{vmatrix}
	- b \begin{vmatrix} d & f \\ g & i \end{vmatrix}
	+ c \begin{vmatrix} d & e \\ g & h \end{vmatrix}.
	$$
	
	Expanding the $2 \times 2$ determinants (minors):
	$$
	\text{det}(A) = a(ei - fh) - b(di - fg) + c(dh - eg).
	$$
	
	For higher dimension matrices, ask a computer (it's outside the scope of this course).
	
	\noindent \textbf{Example:} For the matrix
	$$
	B = \begin{pmatrix} 1 & 2 \\ 2 & 4 \end{pmatrix},
	$$
	$\det{B} = 0$, hence, the matrix is not invertible (you can try to invert it, you will encounter a problem).
	
	While for the matrix
	$$
	A = \begin{pmatrix} 1 & 2 \\ 3 & 4 \end{pmatrix},
	$$
	$\det{A} = -2 \neq 0$, hence, the matrix is invertible (and we found the inverse in the previous section).

	\noindent \textbf{Example ($3 \times 3$):} For the matrix
	$$
	N = \begin{pmatrix} 1 & 2 & 3 \\ 0 & 1 & 4 \\ 5 & 6 & 0 \end{pmatrix}
	$$
	of the previous section, expanding along the first row:
	$$
	\det(N) =
	1 \begin{vmatrix} 1 & 4 \\ 6 & 0 \end{vmatrix}
	- 2 \begin{vmatrix} 0 & 4 \\ 5 & 0 \end{vmatrix}
	+ 3 \begin{vmatrix} 0 & 1 \\ 5 & 6 \end{vmatrix}
	$$
	$$
	= 1 \times (0 - 24) - 2 \times (0 - 20) + 3 \times (0 - 5) = -24 + 40 - 15 = 1 \neq 0
	$$
	so $N$ is invertible, which is consistent with the inverse we computed in the previous section.

	\noindent \textbf{Remark:} the determinant is often the quickest way to know \textit{whether} a matrix can be inverted, but it does not give the inverse.
	When a matrix is \textit{not} invertible, it means its columns are not linearly independent: for $B = \begin{pmatrix} 1 & 2 \\ 2 & 4 \end{pmatrix}$ above, the second column is twice the first.
		
	\section{Eigenvalues}
	
	Let $A$ be an $n \times n$ matrix.
	A scalar $\lambda$ is called an \textbf{eigenvalue} of $A$ if there exists a non-zero vector $\mathbf{v} \in \mathbb{R}^n$ (or $\mathbb{C}^n$) such that
	$$
	A \mathbf{v} = \lambda \mathbf{v}.
	$$
	The vector $\mathbf{v}$ is referred to as the corresponding \textbf{eigenvector} of $A$ associated with $\lambda$.
	
	\paragraph{Finding Eigenvalues}
	To find the eigenvalues of a matrix $ A $, we solve the characteristic equation:
	$$
	\det(A - \lambda I) = 0,
	$$
	where $ I $ is the identity matrix of the same size as $ A $, and $ \det(\cdot) $ denotes the determinant. The solutions $ \lambda $ are the eigenvalues of $ A $.
	
	\noindent \textbf{Example:}
	Consider the matrix
	$$
	A = \begin{pmatrix}
		4 & 1 \\
		2 & 3
	\end{pmatrix}.
	$$
	The characteristic equation is given by:
	$$
	\det(A - \lambda I) = \det\left(\begin{pmatrix}
		4 & 1 \\
		2 & 3
	\end{pmatrix} - \lambda \begin{pmatrix}
		1 & 0 \\
		0 & 1
	\end{pmatrix}\right)
	$$
	$$ = \det\left(\begin{pmatrix}
		4-\lambda & 1 \\
		2 & 3-\lambda
	\end{pmatrix}\right) = (4-\lambda)(3-\lambda) - 2 = \lambda^2 - 7\lambda + 10 = 0.
	$$
	The eigenvalues are the roots of this quadratic equation, $ \lambda_1 = 5 $ and $ \lambda_2 = 2 $.

	\noindent \textbf{Example (triangular matrix):}
	For a triangular matrix, the eigenvalues can be read directly on the diagonal.
	Indeed, for
	$$
	T = \begin{pmatrix} 2 & 7 & -1 \\ 0 & 5 & 3 \\ 0 & 0 & 1 \end{pmatrix},
	$$
	the matrix $T - \lambda I$ is still triangular, and the determinant of a triangular matrix is the product of its diagonal entries:
	$$
	\det(T - \lambda I) = (2-\lambda)(5-\lambda)(1-\lambda) = 0
	\qquad \Longrightarrow \qquad
	\lambda = 2, \ \lambda = 5 \ \text{ or } \ \lambda = 1
	$$

	\noindent \textbf{Example (complex eigenvalues):}
	Eigenvalues are not always real.
	For the rotation matrix
	$$
	R = \begin{pmatrix} 0 & -1 \\ 1 & 0 \end{pmatrix},
	$$
	the characteristic equation is $\det(R - \lambda I) = \lambda^2 + 1 = 0$, whose solutions are $\lambda = i$ and $\lambda = -i$.
	This makes sense geometrically: $R$ rotates every vector of $\mathbb{R}^2$ by $90^\circ$, so no non-zero \textit{real} vector keeps its direction.

	\noindent \textbf{Example (repeated eigenvalue):}
	For
	$$
	S = \begin{pmatrix} 3 & 1 \\ 0 & 3 \end{pmatrix},
	$$
	the characteristic equation is $(3-\lambda)^2 = 0$, so $\lambda = 3$ is a \textit{single} eigenvalue, counted with multiplicity $2$.
	We will see in the section on diagonalization why this case deserves attention.
	
	\section{Eigenvectors}
	
	Given an eigenvalue $\lambda$ of a matrix $A$, the corresponding \textbf{eigenvector} $\mathbf{v}$ is any non-zero vector that satisfies the equation:
	$$
	A \mathbf{v} = \lambda \mathbf{v}.
	$$
	To find the eigenvectors associated with $\lambda$, we solve the system:
	$$
	(A - \lambda I)\mathbf{v} = 0.
	$$
	This is a system of linear equations.
	
	\noindent \textbf{Example:}
	For the matrix $ A $ from the previous example and $ \lambda_1 = 5 $, we solve:
	$$
	(A - 5I)\mathbf{v} = \begin{pmatrix}
		-1 & 1 \\
		2 & -2
	\end{pmatrix} \mathbf{v} = \mathbf{0}.
	$$
	One solution to this system is $\mathbf{v}_1 = \begin{pmatrix} 1 \\ 1 \end{pmatrix}$
	\footnote{Any vector of the form $\mathbf{v} = \begin{pmatrix} \nu \\ \nu \end{pmatrix}, \nu \in \R^*$ would work, it is common practice to set $\nu=1$.}.
	
	Similarly, for $ \lambda_2 = 2 $, we solve:
	$$
	(A - 2I)\mathbf{v} = \begin{pmatrix}
		2 & 1 \\
		2 & 1
	\end{pmatrix} \mathbf{v} = \mathbf{0},
	$$
	which gives $\mathbf{v}_2 = \begin{pmatrix} -1 \\ 2 \end{pmatrix}$ 	\footnote{Again, any vector of the form $\mathbf{v} = \begin{pmatrix} -\nu \\ 2\nu \end{pmatrix}, \nu \in \R^*$ would work, it is common practice to set $\nu=1$.}.

	\noindent \textbf{Example (in practice):}
	Rather than inverting anything, it is often quickest to write $A\mathbf{v} = \lambda \mathbf{v}$ line by line.
	For $A = \begin{pmatrix} 4 & 1 \\ 2 & 3 \end{pmatrix}$ and $\lambda_1 = 5$, with $\mathbf{v} = \begin{pmatrix} x \\ y \end{pmatrix}$:
	$$
	\begin{cases}
		4x + y = 5x \\
		2x + 3y = 5y
	\end{cases}
	\qquad \Longleftrightarrow \qquad
	\begin{cases}
		y = x \\
		2x = 2y
	\end{cases}
	\qquad \Longleftrightarrow \qquad
	y = x
	$$
	The two equations say the same thing, which is expected: if they did not, the only solution would be $\mathbf{v} = \mathbf{0}$, and $\lambda$ would not be an eigenvalue.
	Setting $x = 1$ gives back $\mathbf{v}_1 = \begin{pmatrix} 1 \\ 1 \end{pmatrix}$.

	\noindent \textbf{Example (repeated eigenvalue):}
	For the matrix $S = \begin{pmatrix} 3 & 1 \\ 0 & 3 \end{pmatrix}$ above, with its single eigenvalue $\lambda = 3$:
	$$
	(S - 3I)\mathbf{v} = \begin{pmatrix} 0 & 1 \\ 0 & 0 \end{pmatrix}\begin{pmatrix} x \\ y \end{pmatrix} = \begin{pmatrix} y \\ 0 \end{pmatrix} = \mathbf{0}
	\qquad \Longrightarrow \qquad
	y = 0
	$$
	so every eigenvector is a multiple of $\begin{pmatrix} 1 \\ 0 \end{pmatrix}$.
	Even though the eigenvalue is repeated \textit{twice}, we only get \textit{one} independent eigenvector: this is exactly the situation that prevents diagonalization.
	
	\section{Diagonalization}
	
	A square matrix $A$ is said to be \textbf{diagonalizable} if there exists an invertible matrix $P$ and a diagonal matrix $D$ such that:
	$$
	A = PDP^{-1}.
	$$
	When this is possible, the diagonal elements of $D$ are the eigenvalues of $A$, and the columns of $P$ are the corresponding eigenvectors.
	
	\paragraph{Procedure for Diagonalization}
	
	\begin{enumerate}
		\item Find the eigenvalues $\lambda_1, \lambda_2, \ldots, \lambda_n$ of $A$.
		\item For each eigenvalue $\lambda_i$, find the corresponding eigenvector $\mathbf{v}_i$.
		\item Form the matrix $P$ using the eigenvectors as columns: $P = \begin{pmatrix} \mathbf{v}_1 & \mathbf{v}_2 & \ldots & \mathbf{v}_n \end{pmatrix}$.
		\item The matrix $D$ is a diagonal matrix with eigenvalues on the diagonal:
		$$
		D =
		\begin{pmatrix}
			\lambda_1 & & 0 & 0 & 0 \\
			 & \lambda_2 & & \ddots & 0 \\
			0 & & \ddots & & 0 \\
			0 & \ddots & & \ddots &  \\
			0 & 0 & 0 & & \lambda_n
		\end{pmatrix}.
		$$
		\item Verify that $A = PDP^{-1}$.
	\end{enumerate}
	
	\noindent \textbf{Example:}
	Consider the matrix $A = \begin{pmatrix} 4 & 1 \\ 2 & 3 \end{pmatrix}$.
	From our earlier work, the eigenvalues are $\lambda_1 = 5$ and $\lambda_2 = 2$, with corresponding eigenvectors $\mathbf{v}_1 = \begin{pmatrix} 1 \\ 1 \end{pmatrix}$ and $\mathbf{v}_2 = \begin{pmatrix} -1 \\ 2 \end{pmatrix}$.
	
	Thus, we can form:
	$$
	P = \begin{pmatrix} 1 & -1 \\ 1 & 2 \end{pmatrix},
	\quad D = \begin{pmatrix} 5 & 0 \\ 0 & 2 \end{pmatrix}.
	$$
	Finally, verify that:
	$$
	P^{-1} = \frac{1}{3} \begin{pmatrix} 2 & 1 \\ -1 & 1 \end{pmatrix},
	\quad A = PDP^{-1} = \begin{pmatrix} 4 & 1 \\ 2 & 3 \end{pmatrix}.
	$$
	In detail:
	$$
	PDP^{-1} = \frac{1}{3}\begin{pmatrix} 1 & -1 \\ 1 & 2 \end{pmatrix}\begin{pmatrix} 5 & 0 \\ 0 & 2 \end{pmatrix}\begin{pmatrix} 2 & 1 \\ -1 & 1 \end{pmatrix}
	= \frac{1}{3}\begin{pmatrix} 5 & -2 \\ 5 & 4 \end{pmatrix}\begin{pmatrix} 2 & 1 \\ -1 & 1 \end{pmatrix}
	= \frac{1}{3}\begin{pmatrix} 12 & 3 \\ 6 & 9 \end{pmatrix}
	= \begin{pmatrix} 4 & 1 \\ 2 & 3 \end{pmatrix} \ \checkmark
	$$

	\noindent \textbf{Counter-example (not every matrix is diagonalizable):}
	Consider again $S = \begin{pmatrix} 3 & 1 \\ 0 & 3 \end{pmatrix}$.
	Its only eigenvalue is $\lambda = 3$, and all its eigenvectors are multiples of $\begin{pmatrix} 1 \\ 0 \end{pmatrix}$.
	We can therefore not build a $2 \times 2$ matrix $P$ with two \textit{linearly independent} eigenvectors as columns: any such $P$ would have two proportional columns, hence $\det(P) = 0$, and $P$ would not be invertible.
	The matrix $S$ is \textbf{not diagonalizable}.

	In general, an $n \times n$ matrix is diagonalizable if and only if it has $n$ linearly independent eigenvectors.
	Two useful sufficient conditions:
	\begin{itemize}
		\item if the matrix has $n$ \textit{distinct} eigenvalues, it is diagonalizable (as in the example above);
		\item if the matrix is \textit{symmetric} (with real entries), it is diagonalizable, and $P$ can even be chosen orthogonal. This is the case that matters for Principal Component Analysis.
	\end{itemize}

	\noindent \textbf{Why bother?}
	Because a diagonal matrix is much easier to compute with.
	For instance, powers of $A$ become trivial: since
	$$
	A^k = \underbrace{PDP^{-1} \cdot PDP^{-1} \cdots PDP^{-1}}_{k \text{ times}} = P D^k P^{-1}
	$$
	(the inner $P^{-1}P = I$ cancel out), and $D^k$ is simply the diagonal matrix of the $\lambda_i^k$, we get for our example
	$$
	A^{k} = \frac{1}{3}\begin{pmatrix} 1 & -1 \\ 1 & 2 \end{pmatrix}\begin{pmatrix} 5^{k} & 0 \\ 0 & 2^{k} \end{pmatrix}\begin{pmatrix} 2 & 1 \\ -1 & 1 \end{pmatrix}
	$$
	without ever multiplying $A$ by itself.
	This is what makes diagonalization useful for iterated processes (Markov chains, dynamical systems, \dots).
	
	
	\section{Norms}
	In vector spaces, a \textbf{norm} is a function that assigns a non-negative length or size to vectors.
	Norms are widely used to measure the magnitude of vectors.
	The most common norms are the $L_1$, $L_2$, and $L_p$ norms.
	
	\paragraph{$L_1$ Norm (Manhattan Norm)}
	The $L_1$ norm of a vector $\mathbf{v} = (v_1, v_2, \dots, v_n)$ in $\mathbb{R}^n$ is defined as the sum of the absolute values of its components:
	$$
	\|\mathbf{v}\|_1 = |v_1| + |v_2| + \dots + |v_n| = \sum_{i=1}^{n} |v_i|.
	$$
	This norm is also known as the \textbf{Manhattan norm} or \textbf{taxicab norm}, as it represents the distance between two points in a grid-based path, like the blocks in a city.
	
	\noindent
	\textbf{Example:}
	For the vector $\mathbf{v} = \begin{pmatrix} 3 \\ -4 \\ 1 \end{pmatrix}$ in $\mathbb{R}^3$, the $L_1$ norm is:
	$$
	\|\mathbf{v}\|_1 = |3| + |-4| + |1| = 3 + 4 + 1 = 8.
	$$
	
	\paragraph{$L_2$ Norm (Euclidean Norm)}
	The $L_2$ norm of a vector $\mathbf{v} = (v_1, v_2, \dots, v_n)$ is the square root of the sum of the squares of its components:
	$$
	\|\mathbf{v}\|_2 = \sqrt{v_1^2 + v_2^2 + \dots + v_n^2} = \sqrt{\sum_{i=1}^{n} v_i^2}.
	$$
	This norm is commonly known as the \textbf{Euclidean norm} because it represents the usual notion of distance in Euclidean space.
	
	\noindent
	\textbf{Example:}
	For the vector $\mathbf{v} = \begin{pmatrix} 3 \\ -4 \\ 1 \end{pmatrix}$ in $\mathbb{R}^3$, the $L_2$ norm is:
	$$
	\|\mathbf{v}\|_2 = \sqrt{3^2 + (-4)^2 + 1^2} = \sqrt{9 + 16 + 1} = \sqrt{26}.
	$$
	
	\paragraph{$L_p$ Norm}
	The $L_p$ norm generalizes the $L_1$ and $L_2$ norms and is defined for any real number $p \geq 1$.
	The $L_p$ norm of a vector $\mathbf{v} = (v_1, v_2, \dots, v_n)$ is given by:
	$$
	\|\mathbf{v}\|_p = \left( |v_1|^p + |v_2|^p + \dots + |v_n|^p \right)^{\frac{1}{p}} = \left(\sum_{i=1}^{n} |v_i|^p\right)^{\frac{1}{p}}.
	$$
	Different values of $p$ result in different norms.
	The $L_2$ norm is a special case where $p = 2$, and the $L_1$ norm is a special case where $p = 1$.
	
	\noindent
	\textbf{Example:}
	For the vector $\mathbf{v} = \begin{pmatrix} 3 \\ -4 \\ 1 \end{pmatrix}$ in $\mathbb{R}^3$, the $L_p$ norm with $p = 3$ is:
	$$
	\|\mathbf{v}\|_3 = \left( |3|^3 + |-4|^3 + |1|^3 \right)^{\frac{1}{3}} = \left( 27 + 64 + 1 \right)^{\frac{1}{3}} = \sqrt[3]{92}.
	$$
	
	\paragraph{Properties of Norms}
	To be a norm, a function $\| \cdot \|: \ \R^n \to \R$ must verify:
	\begin{enumerate}
		\item \textbf{Non-negativity:}
		$\|\mathbf{v}\| \geq 0$ for all vectors $\mathbf{v}$, and $\|\mathbf{v}\| = 0$ if and only if $\mathbf{v} = \mathbf{0}$.
		\item \textbf{Scalar Multiplication:}
		For any scalar $\alpha$ and vector $\mathbf{v}$, $\|\alpha \mathbf{v}\| = |\alpha| \|\mathbf{v}\|$.
		\item \textbf{Triangle Inequality:}
		For any vectors $\mathbf{v}$ and $\mathbf{w}$, $\|\mathbf{v} + \mathbf{w}\| \leq \|\mathbf{v}\| + \|\mathbf{w}\|$.
	\end{enumerate}

	\noindent
	\textbf{Example:}
	Take $\mathbf{v} = \begin{pmatrix} 3 \\ -4 \\ 1 \end{pmatrix}$ and $\mathbf{w} = \begin{pmatrix} -1 \\ 2 \\ 2 \end{pmatrix}$, so that $\mathbf{v} + \mathbf{w} = \begin{pmatrix} 2 \\ -2 \\ 3 \end{pmatrix}$.
	\begin{itemize}
		\item \textit{Scalar multiplication:} $2\mathbf{v} = \begin{pmatrix} 6 \\ -8 \\ 2 \end{pmatrix}$, and $\|2\mathbf{v}\|_1 = 6 + 8 + 2 = 16 = 2 \times 8 = |2| \, \|\mathbf{v}\|_1$. \checkmark
		\item \textit{Triangle inequality, for the $L_1$ norm:}
		$$
		\|\mathbf{v} + \mathbf{w}\|_1 = 2 + 2 + 3 = 7
		\qquad \leq \qquad
		\|\mathbf{v}\|_1 + \|\mathbf{w}\|_1 = 8 + 5 = 13 \ \checkmark
		$$
		\item \textit{Triangle inequality, for the $L_2$ norm:}
		$$
		\|\mathbf{v} + \mathbf{w}\|_2 = \sqrt{4+4+9} = \sqrt{17} \approx 4.12
		\qquad \leq \qquad
		\|\mathbf{v}\|_2 + \|\mathbf{w}\|_2 = \sqrt{26} + 3 \approx 8.10 \ \checkmark
		$$
	\end{itemize}
	The inequality is an equality only when $\mathbf{v}$ and $\mathbf{w}$ point in the same direction; here they do not, so there is some slack.

	\noindent
	\textbf{Why it matters in practice:}
	when the vector measures an \textit{error}, the choice of norm changes what the model is asked to avoid.
	Compare the two error vectors
	$$
	\mathbf{e}_1 = \begin{pmatrix} 1 \\ 1 \\ 1 \\ 1 \end{pmatrix}
	\qquad \text{and} \qquad
	\mathbf{e}_2 = \begin{pmatrix} 4 \\ 0 \\ 0 \\ 0 \end{pmatrix}
	$$
	They have the \textit{same} $L_1$ norm ($\|\mathbf{e}_1\|_1 = \|\mathbf{e}_2\|_1 = 4$), but different $L_2$ norms: $\|\mathbf{e}_1\|_2 = 2$ while $\|\mathbf{e}_2\|_2 = 4$.
	An $L_2$ loss strongly prefers many small errors over one large one; an $L_1$ loss is indifferent between the two.
	
	
	
	\section{The Gram-Schmidt Algorithm}	
	We assume the following:
	\begin{itemize}
		\item The vectors $\{ \mathbf{v}_1, \mathbf{v}_2, \dots, \mathbf{v}_n \}$ are linearly independent in an inner product space $V$.
		\item The space $V$ is equipped with an inner product $\langle \mathbf{u}, \mathbf{v} \rangle$ which is usually the standard dot product in $\mathbb{R}^n$.
	\end{itemize}
	
	Given a set of linearly independent vectors $\mathbf{v}_1, \mathbf{v}_2, \dots, \mathbf{v}_n$, the goal is to find an orthogonal set of vectors $\mathbf{u}_1, \mathbf{u}_2, \dots, \mathbf{u}_n$ such that:
	\begin{itemize}
		\item $\text{span}(\mathbf{v}_1, \dots, \mathbf{v}_k) = \text{span}(\mathbf{u}_1, \dots, \mathbf{u}_k)$ for all $k$,
		\item $\langle \mathbf{u}_i, \mathbf{u}_j \rangle = 0$ for all $i \neq j$ (orthogonality),
		\item Optionally, we can normalize the vectors so that $\|\mathbf{u}_i\| = 1$, forming an orthonormal set.
	\end{itemize}
	Let $\mathbf{v}_1, \mathbf{v}_2, \dots, \mathbf{v}_n$ be a set of linearly independent vectors.
	The Gram-Schmidt process proceeds as follows:
	
	\paragraph{Step 1: Construct the first orthogonal vector}
	\noindent\\
	We begin by setting the first vector in the orthogonal set equal to the first vector in the original set:
	$$ \mathbf{u}_1 = \mathbf{v}_1 $$
	At this point, $\mathbf{u}_1$ is trivially orthogonal because it's the only vector in the set.
	
	\paragraph{Step 2: Orthogonalize the remaining vectors}
	\noindent\\
	For each subsequent vector $\mathbf{v}_k$ (where $k = 2, 3, \dots, n$), we subtract off the projections onto the previously calculated orthogonal vectors $\mathbf{u}_1, \dots, \mathbf{u}_{k-1}$.
	
	The orthogonal vector $\mathbf{u}_k$ is computed as:
	$$
	\mathbf{u}_k = \mathbf{v}_k - \sum_{j=1}^{k-1} \frac{\langle \mathbf{v}_k, \mathbf{u}_j \rangle}{\langle \mathbf{u}_j, \mathbf{u}_j \rangle} \mathbf{u}_j
	$$
	
	This formula ensures that $\mathbf{u}_k$ is orthogonal to each of the previous vectors $\mathbf{u}_1, \dots, \mathbf{u}_{k-1}$.
	
	\paragraph{Step 3: Normalize (Optional)}
	\noindent\\
	If we want an orthonormal set of vectors, we normalize each $\mathbf{u}_k$ by dividing it by its magnitude:
	$$
	\mathbf{e}_k = \frac{\mathbf{u}_k}{\|\mathbf{u}_k\|}
	$$
	The vectors $\mathbf{e}_1, \dots, \mathbf{e}_n$ now form an orthonormal basis for the same subspace.

	\paragraph{Worked Example}
	\noindent\\
	Let us orthogonalize the following linearly independent vectors of $\mathbb{R}^3$:
	$$
	\mathbf{v}_1 = \begin{pmatrix} 1 \\ 1 \\ 0 \end{pmatrix}, \quad
	\mathbf{v}_2 = \begin{pmatrix} 1 \\ 0 \\ 1 \end{pmatrix}, \quad
	\mathbf{v}_3 = \begin{pmatrix} 0 \\ 1 \\ 1 \end{pmatrix}
	$$

	\noindent \textbf{Step 1.} We simply take
	$$
	\mathbf{u}_1 = \mathbf{v}_1 = \begin{pmatrix} 1 \\ 1 \\ 0 \end{pmatrix},
	\qquad \text{with} \qquad
	\langle \mathbf{u}_1, \mathbf{u}_1 \rangle = 1 + 1 + 0 = 2
	$$

	\noindent \textbf{Step 2 (second vector).} Since $\langle \mathbf{v}_2, \mathbf{u}_1 \rangle = 1 \times 1 + 0 \times 1 + 1 \times 0 = 1$,
	$$
	\mathbf{u}_2 = \mathbf{v}_2 - \frac{\langle \mathbf{v}_2, \mathbf{u}_1 \rangle}{\langle \mathbf{u}_1, \mathbf{u}_1 \rangle} \mathbf{u}_1
	= \begin{pmatrix} 1 \\ 0 \\ 1 \end{pmatrix} - \frac{1}{2}\begin{pmatrix} 1 \\ 1 \\ 0 \end{pmatrix}
	= \begin{pmatrix} \frac{1}{2} \\ -\frac{1}{2} \\ 1 \end{pmatrix}
	$$
	We check that it worked: $\langle \mathbf{u}_2, \mathbf{u}_1 \rangle = \frac{1}{2} - \frac{1}{2} + 0 = 0$. \checkmark

	\noindent \textbf{Step 2 (third vector).} We now subtract \textit{two} projections.
	With $\langle \mathbf{v}_3, \mathbf{u}_1 \rangle = 0 + 1 + 0 = 1$, $\langle \mathbf{v}_3, \mathbf{u}_2 \rangle = 0 - \frac{1}{2} + 1 = \frac{1}{2}$ and $\langle \mathbf{u}_2, \mathbf{u}_2 \rangle = \frac{1}{4} + \frac{1}{4} + 1 = \frac{3}{2}$:
	$$
	\mathbf{u}_3 = \mathbf{v}_3
	- \frac{\langle \mathbf{v}_3, \mathbf{u}_1 \rangle}{\langle \mathbf{u}_1, \mathbf{u}_1 \rangle} \mathbf{u}_1
	- \frac{\langle \mathbf{v}_3, \mathbf{u}_2 \rangle}{\langle \mathbf{u}_2, \mathbf{u}_2 \rangle} \mathbf{u}_2
	= \begin{pmatrix} 0 \\ 1 \\ 1 \end{pmatrix}
	- \frac{1}{2}\begin{pmatrix} 1 \\ 1 \\ 0 \end{pmatrix}
	- \frac{1/2}{3/2}\begin{pmatrix} \frac{1}{2} \\ -\frac{1}{2} \\ 1 \end{pmatrix}
	$$
	$$
	= \begin{pmatrix} 0 - \frac{1}{2} - \frac{1}{6} \\ 1 - \frac{1}{2} + \frac{1}{6} \\ 1 - \frac{1}{3} \end{pmatrix}
	= \begin{pmatrix} -\frac{2}{3} \\ \frac{2}{3} \\ \frac{2}{3} \end{pmatrix}
	$$
	Again we check: $\langle \mathbf{u}_3, \mathbf{u}_1 \rangle = -\frac{2}{3} + \frac{2}{3} + 0 = 0$ \checkmark \ and $\langle \mathbf{u}_3, \mathbf{u}_2 \rangle = -\frac{1}{3} - \frac{1}{3} + \frac{2}{3} = 0$. \checkmark

	\noindent \textbf{Step 3 (normalization).} The norms are
	$$
	\|\mathbf{u}_1\| = \sqrt{2}, \qquad
	\|\mathbf{u}_2\| = \sqrt{\frac{3}{2}} = \frac{\sqrt{6}}{2}, \qquad
	\|\mathbf{u}_3\| = \sqrt{\frac{4}{3}} = \frac{2}{\sqrt{3}}
	$$
	so the orthonormal basis is
	$$
	\mathbf{e}_1 = \frac{1}{\sqrt{2}}\begin{pmatrix} 1 \\ 1 \\ 0 \end{pmatrix}, \qquad
	\mathbf{e}_2 = \frac{1}{\sqrt{6}}\begin{pmatrix} 1 \\ -1 \\ 2 \end{pmatrix}, \qquad
	\mathbf{e}_3 = \frac{1}{\sqrt{3}}\begin{pmatrix} -1 \\ 1 \\ 1 \end{pmatrix}
	$$
	(for $\mathbf{e}_2$ and $\mathbf{e}_3$, we first factored out $\frac{1}{2}$ and $\frac{2}{3}$ respectively, which is allowed since scaling a vector does not change its direction).

	\noindent \textbf{Remark:} the order of the vectors matters.
	Starting from $\mathbf{v}_2$ instead of $\mathbf{v}_1$ would give a different orthonormal basis, spanning the same space.
	Only the first vector is kept unchanged (up to its length) by the algorithm.
	
	
	
\end{document}
